\section{A signed-kernel representation}
\label{signed:sec:kernel}

\begin{definition}\label{signed:def:kappa}
For real $b>0$ and $T>0$, define
\begin{equation}\label{signed:eq:kappa-base}
 \kappa_{1,b}(e^{-T})=
 -\frac{T^b}{\Gamma(b+1)}
 +\frac1{\Gamma(b)}\int_T^\infty\frac{v^{b-1}}{e^v-1}\dd v.
\end{equation}
For an integer $a\ge2$, define successively
\begin{equation}\label{signed:eq:kappa-recursion}
 \kappa_{a,b}(u)=\int_u^1\frac{\kappa_{a-1,b}(v)}{v}\dd v,
 \qquad 0<u<1.
\end{equation}
\end{definition}

\begin{theorem}[Moments and one sign change]\label{signed:thm:kernel}
For every integer $a\ge1$ and real $b>0$, the density $\kappa_{a,b}$ is integrable on $(0,1)$ and
\begin{equation}\label{signed:eq:kernel-moments}
 \int_0^1u^{n-1}\kappa_{a,b}(u)\dd u
 =\frac{H_{n-1}^{(b)}}{n^a},\qquad n\ge1.
\end{equation}
In particular its total mass is zero. There is exactly one $r_{a,b}\in(0,1)$ such that
\begin{equation}\label{signed:eq:kernel-sign}
 \kappa_{a,b}(u)<0\quad(0<u<r_{a,b}),\qquad
 \kappa_{a,b}(u)>0\quad(r_{a,b}<u<1).
\end{equation}
The crossing is strict, and $r_{a,b}<r_{a-1,b}$ when $a\ge2$.
\end{theorem}
\begin{proof}
First consider $a=1$. The negative term in \eqref{signed:eq:kappa-base} has moment $-n^{-b-1}$. Substituting $u=e^{-T}$ and interchanging nonnegative integrands in the positive term gives
\begin{align*}
 &\frac1{\Gamma(b)}\int_0^\infty e^{-nT}
       \int_T^\infty\frac{v^{b-1}}{e^v-1}\dd v\dd T\\
 &\hspace{12mm}=\frac1{n\Gamma(b)}\int_0^\infty
       \frac{v^{b-1}(1-e^{-nv})}{e^v-1}\dd v
 =\frac1n\sum_{k=1}^n k^{-b}.
\end{align*}
The final equality uses the finite geometric identity
$(1-e^{-nv})/(e^v-1)=\sum_{k=1}^n e^{-kv}$.
Subtracting $n^{-b-1}$ proves the base moment formula. At $n=1$, both positive and negative terms have finite integral~1, so integrability and zero mass are established without subtracting divergent quantities.

As a function of $T$, the base density has derivative
\begin{equation}\label{signed:eq:kappa-base-derivative}
 \frac{\mathrm d}{\mathrm dT}\kappa_{1,b}(e^{-T})
 =-\frac{T^{b-1}}{\Gamma(b)(1-e^{-T})}<0.
\end{equation}
It is positive as $T\downarrow0$: the limit is $\zeta(b)>0$ if $b>1$, and is $+\infty$ if $0<b\le1$. It tends to $-\infty$ as $T\to\infty$. Thus it has exactly one crossing, with the sign pattern stated in the $u$ coordinate.

For the induction, the integral operator in \eqref{signed:eq:kappa-recursion} is an $L^1$ contraction:
\[
 \int_0^1\left|\int_u^1\frac{f(v)}v\dd v\right|\dd u
 \le\int_0^1|f(v)|\dd v.
\]
Fubini's theorem gives the moment of the new density as $1/n$ times the previous moment; in particular it preserves zero mass. If the preceding crossing is $r$, the new density is positive on $[r,1)$, while on $(0,r)$ its derivative is $-\kappa_{a-1,b}(u)/u>0$. Zero mass forces it to be negative somewhere on that interval. Strict increase then gives exactly one crossing in $(0,r)$. This completes the induction.
\end{proof}

\begin{corollary}[Integral continuation]\label{signed:cor:stieltjes}
For $z\in\mathbb C\setminus[1,\infty)$,
\begin{equation}\label{signed:eq:stieltjes}
 F_{a,b}(z)=\int_0^1\frac{z}{1-zu}\kappa_{a,b}(u)\dd u.
\end{equation}
In particular, if $0<\theta<\pi$,
\begin{equation}\label{signed:eq:angular-integral}
 \Imag F_{a,b}(e^{\iu\theta})
 =\sin\theta\,J_{a,b}(\cos\theta),\qquad
 J_{a,b}(c)=\int_0^1\frac{\kappa_{a,b}(u)}{1-2cu+u^2}\dd u.
\end{equation}
\end{corollary}
\begin{proof}
For $|z|<1$, expand $z/(1-zu)$ and use the moments. On compact subsets of the slit plane the denominator is bounded away from zero uniformly in $u\in[0,1]$, so the integral is analytic there. This continues the disk identity. Taking the imaginary part of $e^{\iu\theta}/(1-ue^{\iu\theta})$ gives \eqref{signed:eq:angular-integral}.
\end{proof}

\begin{lemma}[A sign test for zero-mass densities]\label{signed:lem:sign-test}
Let $\kappa$ have zero integral and the strict sign pattern \eqref{signed:eq:kernel-sign}, with crossing $r$. If $q$ is strictly decreasing and the integral exists, then $\int q\kappa<0$. If $q$ is strictly increasing, the sign is positive.
\end{lemma}
\begin{proof}
Subtract $q(r)\int\kappa=0$. The product $(q(u)-q(r))\kappa(u)$ has the asserted strict sign on each side of $r$.
\end{proof}

\section{The unique zero on the upper semicircle}
\label{signed:sec:zeros}

\begin{theorem}[Existence, uniqueness, and simplicity]\label{signed:thm:unique-zero}
Let $a\ge1$ be an integer and let $b>0$ be real. There is exactly one
$\theta_{a,b}\in(0,\pi)$ such that
$\Imag F_{a,b}(e^{\iu\theta_{a,b}})=0$. It satisfies
\begin{equation}\label{signed:eq:zero-range}
 0<\theta_{a,b}<\frac\pi2.
\end{equation}
The imaginary part is positive for $0<\theta<\theta_{a,b}$ and negative for $\theta_{a,b}<\theta<\pi$. At its zero the angular derivative is strictly negative. Consequently,
\begin{equation}\label{signed:eq:g-negative}
 g_{a,b}<0.
\end{equation}
\end{theorem}
\begin{proof}
Suppress the indices in $J$ and $\kappa$. The kernel $u\mapsto(1+u^2)^{-1}$ is strictly decreasing, so Lemma~\ref{signed:lem:sign-test} gives $J(0)<0$.

We next show that the limit of $J(c)$ as $c\uparrow1$ is positive, possibly infinite. The negative density is supported in $(0,r)$, away from the singular endpoint, and its contribution has a finite limit. On $(r,1)$ monotone convergence applies. Expanding $(1-u)^{-2}$, treating these two parts separately, and using \eqref{signed:eq:kernel-moments} yields
\begin{equation}\label{signed:eq:J-at-one}
 \lim_{c\uparrow1}J(c)
 =\int_0^1\frac{\kappa(u)}{(1-u)^2}\dd u
 =\sum_{n\ge0}\frac{H_n^{(b)}}{(n+1)^{a-1}}>0.
\end{equation}
Every summand from $n=1$ onward is positive. Thus continuity produces at least one zero $c_0\in(0,1)$.

For uniqueness, write $K_c(u)=(1-2cu+u^2)^{-1}$. If $c'>c$, then
\begin{equation}\label{signed:eq:kernel-ratio}
 \frac{\mathrm d}{\mathrm du}\frac{K_{c'}(u)}{K_c(u)}
 =\frac{2(c'-c)(1-u^2)}{(1-2c'u+u^2)^2}>0.
\end{equation}
At a zero $J(c)=0$, the signed density $K_c\kappa$ still has one sign change and zero mass. Applying Lemma~\ref{signed:lem:sign-test} to the ratio in \eqref{signed:eq:kernel-ratio} gives $J(c')>0$ for every $c'>c$. The same argument with a smaller parameter gives $J(c')<0$ for $c'<c$. There can be only one zero in $(-1,1)$, with the indicated signs.

Finally,
\[
 J'(c_0)=\int_0^1 K_{c_0}(u)\kappa(u)
        \frac{2u}{1-2c_0u+u^2}\dd u>0,
\]
because the last factor has derivative $2(1-u^2)/(1-2c_0u+u^2)^2>0$. Equation~\eqref{signed:eq:angular-integral} therefore has derivative
$-\sin^2\theta_{a,b}\,J'(c_0)<0$ at $\theta_{a,b}=\arccos c_0$. Since $c_0>0$, this zero is below $\pi/2$. The Gaussian inequality is its value at $\theta=\pi/2$.
\end{proof}

\begin{remark}[The first zero is elementary]
The identity $F_{1,1}(z)=\tfrac12\log^2(1-z)$ gives
\[
 \Imag F_{1,1}(e^{\iu\theta})
 =\log\!\left(2\sin\frac\theta2\right)\frac{\theta-\pi}{2}.
\]
Hence $\theta_{1,1}=\pi/3$ and $g_{1,1}=-\pi\log2/8$. This example also checks the direction of the sign change in Theorem~\ref{signed:thm:unique-zero}.
\end{remark}

\section{A four-term asymptotic for the zero}
\label{signed:sec:zero-asymptotic}
Put $\delta_{a,b}=\pi/2-\theta_{a,b}$. The first surviving Fourier mode at $\pi/2$ is the $n=3$ term, whereas the $n=2$ term controls the derivative. Their ratio suggests $(2/3)^a$; the theorem below identifies both the next Fourier corrections and the first nonlinear correction.

\begin{theorem}[Uniform zero asymptotic]\label{signed:thm:zero-asymptotic}
For integer $a\to\infty$, uniformly for real $b\ge1$, set
$A=H_2^{(b)}$, $B=H_3^{(b)}$, and $C=H_4^{(b)}$. Then
\begin{align}\label{signed:eq:zero-asymptotic}
 \delta_{a,b}
 &=\frac A2\left(\frac23\right)^a
   -\frac C2\left(\frac25\right)^a
   +AB\left(\frac13\right)^a
   -\frac{23A^3}{48}\left(\frac8{27}\right)^a
   +O\!\left(\left(\frac27\right)^a\right).
\end{align}
The implied constant and the threshold for $a$ can be chosen independently of $b\ge1$.
\end{theorem}
\begin{proof}
For sufficiently large $a$, the boundary series and its derivative converge absolutely. Normalize the imaginary part by multiplying by $2^a$, and write $\theta=\pi/2-\delta$. The modes $n=2,3,4,5$ give
\begin{align}\label{signed:eq:zero-fourier}
 0={}&\sin(2\delta)-Aq^a\cos(3\delta)-Bs^a\sin(4\delta)
       +Cr^a\cos(5\delta)\notag\\
    &\quad+O\!\left(\delta\,3^{-a}+(2/7)^a\right),
 \qquad q=\frac23,\quad s=\frac12,\quad r=\frac25.
\end{align}
The $n=6$ mode is $O(\delta3^{-a})$. For $n\ge7$, the bound
$H_{n-1}^{(b)}\le1+\log n$ gives a uniform tail estimate: for $a\ge3$,
\[
 \sum_{n\ge7}(1+\log n)(2/n)^a
 \le (2/7)^a\sum_{n\ge7}(1+\log n)(7/n)^3.
\]
The constant is finite and independent of $b$.

Since $1\le A\le3/2$ and $B,C$ are bounded uniformly in $b$, the normalized expression is negative at $\delta=0$ and positive at $\delta=q^a$ for all sufficiently large $a$. Theorem~\ref{signed:thm:unique-zero} identifies the root in that interval. Thus $\delta=O(q^a)$ uniformly.

Taylor expansion of \eqref{signed:eq:zero-fourier} now gives
\begin{align}\label{signed:eq:zero-expanded}
 0={}&2\delta-Aq^a-4Bs^a\delta+Cr^a
       +\frac92Aq^a\delta^2-\frac43\delta^3
       +O((2/7)^a).
\end{align}
For example, the omitted terms have exponential bases bounded by
$q^5$, $sq^3$, $rq^2$, and $q/3$, all smaller than $2/7$. Insert first
$\delta_0=Aq^a/2$. The linear perturbation $-4Bs^a\delta_0$ requires the correction $AB(qs)^a$, while the $n=5$ mode requires $-Cr^a/2$.
The cubic contribution at $\delta_0$ is
\[
 \frac92Aq^a\delta_0^2-\frac43\delta_0^3
 =\frac{23}{24}A^3q^{3a},
\]
so its compensating correction is $-23A^3q^{3a}/48$.
Substitution of the full displayed approximation in \eqref{signed:eq:zero-expanded} leaves a residual $O((2/7)^a)$; all cross-products of the corrections have smaller exponential bases. The derivative of the normalized expression with respect to $\delta$ is $2+o(1)$, uniformly on this interval. The mean-value theorem therefore converts the residual estimate into the error in \eqref{signed:eq:zero-asymptotic}.
\end{proof}

\begin{center}
\small
\begin{tabular}{@{}rrrr@{}}
\toprule
$a$ & $b$ & $\theta_{a,b}$, approximately & $\delta_{a,b}$, approximately\\
\midrule
8&1&1.541868140889&0.02892818590626\\
8&2&1.546656477391&0.02413984940401\\
8&4&1.550268639150&0.02052768764464\\
12&1&1.565028587360&0.005767739435093\\
12&2&1.565988233290&0.004808093504888\\
12&4&1.566708834682&0.004087492112621\\
20&1&1.570570790993&0.0002255358018452\\
20&2&1.570608378737&0.0001879480579893\\
20&4&1.570636570306&0.0001597564888279\\
\bottomrule
\end{tabular}
\end{center}
These entries are numerical illustrations from a 300-term Fourier sum, not certified root intervals. Existence, uniqueness, and the asymptotic estimate do not depend on the table. At $a=20,40,80$ and $b=1,2,8$, the computed next normalized remainder approaches $H_6^{(b)}/2$, as suggested by the $n=7$ mode. A systematic expansion beyond the proved four terms is a separate research problem.

\section{Certified Euler evaluation}
\label{signed:sec:euler}

\subsection{A one-sided rational enclosure}
Define the difference operator with the sign convention
\[
 (Df)(n)=f(n)-f(n+1).
\]
This is the negative of the usual forward difference; it makes the Euler transformation particularly simple. Let
\begin{equation}\label{signed:eq:euler-data}
 f_{a,b}(n)=\frac{H_{2n}^{(b)}}{(2n+1)^a},\qquad
 d_k=D^kf_{a,b}(0),\qquad
 E_N(a,b)=\sum_{k=0}^{N-1}\frac{d_k}{2^{k+1}}.
\end{equation}

\begin{theorem}[Signed Euler certificate]\label{signed:thm:euler-certificate}
For positive integers $a,b$,
\begin{equation}\label{signed:eq:euler-signs}
 d_0=0,\qquad d_k<0\quad(k\ge1),\qquad
 g_{a,b}=\sum_{k\ge0}\frac{d_k}{2^{k+1}}.
\end{equation}
For every $N\ge1$,
\begin{equation}\label{signed:eq:euler-enclosure}
 E_N(a,b)-C_b2^{-N}\le g_{a,b}\le E_N(a,b),\qquad
 C_b=\begin{cases}1,&b=1,\\ b/(b-1),&b\ge2.\end{cases}
\end{equation}
The approximants decrease strictly from $E_1=0$ to $g_{a,b}$.
\end{theorem}
\begin{proof}
The moment representation gives
\begin{equation}\label{signed:eq:d-kernel}
 d_k=\int_0^1(1-u^2)^k\kappa_{a,b}(u)\dd u.
\end{equation}
For $k=0$ this is zero mass. For $k\ge1$ the weight is strictly decreasing, so Lemma~\ref{signed:lem:sign-test} gives $d_k<0$.

We bound either signed mass of the density. When $b=1$,
\[
 \kappa_{1,1}(u)=\log\frac{u}{1-u},\qquad
 \int_0^1(\kappa_{1,1})_+\dd u=\log2<1.
\]
When $b\ge2$, the positive term of \eqref{signed:eq:kappa-base} is at most $\zeta(b)$, so the positive mass is at most $\zeta(b)<1+1/(b-1)=C_b$. Zero mass makes the negative mass equal to the positive mass. The $L^1$ contraction used in Theorem~\ref{signed:thm:kernel}, together with preservation of zero mass, gives the same bound for every $a$. Since $0\le(1-u^2)^k\le1$, we obtain $|d_k|\le C_b$.

The series in \eqref{signed:eq:euler-signs} is now absolutely convergent. Summing its geometric kernel gives
\[
 \sum_{k\ge0}\frac{d_k}{2^{k+1}}
 =\int_0^1\frac{\kappa_{a,b}(u)}{1+u^2}\dd u
 =g_{a,b},
\]
where the last equality is \eqref{signed:eq:stieltjes} at $z=\iu$. All tail terms are negative, and their total absolute value is at most $C_b\sum_{k\ge N}2^{-k-1}=C_b2^{-N}$. This proves the certificate and strict monotonicity.
\end{proof}

The theorem is stronger than applying a general-purpose extrapolation rule and observing stability. Every endpoint in \eqref{signed:eq:euler-enclosure} is rational. For $N=400$, the interval width is at most $2\cdot2^{-400}<8\cdot10^{-121}$. The accompanying files round the enclosures outward to convenient decimal strings; that display rounding is separate from the internal rational width.

\subsection{An exact linear-work finite formula}
Computing a full triangular difference table uses quadratic work. A finite rearrangement avoids it.

\begin{proposition}[Binomial-tail weights]\label{signed:prop:euler-weights}
For $N\ge1$,
\begin{equation}\label{signed:eq:euler-weighted}
 E_N(a,b)=2^{-N}\sum_{n=0}^{N-1}(-1)^n
 \frac{H_{2n}^{(b)}}{(2n+1)^a}\sum_{j=n+1}^{N}\binom Nj.
\end{equation}
At fixed $a,b$, the right side can be formed with $O(N)$ exact rational or integer updates after choosing $N$. All denominators divide
\begin{equation}\label{signed:eq:denominator-bound}
 2^N\operatorname{lcm}(1,2,\ldots,2N)^{a+b}.
\end{equation}
In particular, the number of bit operations is polynomial in the requested precision at fixed indices.
\end{proposition}
\begin{proof}
Expand $D^kf(0)=\sum_{n=0}^k(-1)^n\binom kn f(n)$ and interchange the two finite sums. The required weight identity is
\[
 \sum_{k=n}^{N-1}2^{N-k-1}\binom kn
 =\sum_{j=n+1}^N\binom Nj.
\]
It follows by induction on $N$ using Pascal's recurrence, or by comparing binomial-tail generating functions. This proves \eqref{signed:eq:euler-weighted}.

The binomial coefficients are generated successively by
$\binom Nn=\binom N{n-1}(N-n+1)/n$, and the cumulative tail is updated by subtracting one coefficient. Each harmonic number is updated by two rational summands. Thus only linearly many updates are needed. Each summand's denominator divides the indicated least-common-multiple power, and the additional factor is $2^N$.
Finally, the elementary bound $\operatorname{lcm}(1,\ldots,2N)\le(2N)!$ gives bit length $O(N+(a+b)N\log N)$ for a common denominator; the numerator has a comparable polynomial bound. Since $N=O(d)$ suffices for $d$ decimal places by \eqref{signed:eq:euler-enclosure}, even elementary integer arithmetic proves polynomial bit complexity. No optimal multiplication or sharp least-common-multiple estimate is required.
\end{proof}

\subsection{Independent depth-one interval arithmetic}
The right sides of the reductions are evaluated using rational enclosures separate from the Gaussian evaluator. For $s\ge1$, the sequences $(2n+1)^{-s}$ and $(n+1)^{-s}$ are moments of positive densities, so all their $D$-differences are nonnegative. Their Euler sums give lower bounds for $\beta(s)$ and $\eta(s)$, with tails at most $2^{-N}$. For integer $s\ge2$, divide the latter interval by $1-2^{1-s}$ to enclose $\zeta(s)$.

For $\pi$, the implementation uses Machin's formula
\[
 \pi=16\arctan(1/5)-4\arctan(1/239)
\]
and the alternating-series bounds for each rational arctangent series. The identity follows directly by tangent addition and the fact that the angle lies in the first quadrant. For $\log2$, it uses
\begin{equation}\label{signed:eq:log2-bound}
 0<\log2-2\sum_{n=0}^{N-1}\frac1{(2n+1)3^{2n+1}}
 \le\frac9{4(2N+1)3^{2N+1}}.
\end{equation}
The expansion is $2\operatorname{arctanh}(1/3)$; the bound replaces the remaining denominators by $2N+1$ and sums a geometric series.

Ordinary rational interval addition and multiplication then enclose every polynomial right side. An overlap check can detect an erroneous formula or implementation if the intervals separate. Overlap alone does not prove an identity; the proofs in Sections~\ref{gaussian:sec:double-parity} and~\ref{gaussian:sec:shuffle} do that.

\subsection{Affine harmonic sequences as signed moments}
\label{affine:sec:moments}
The signed-measure formulation also certifies sequences whose harmonic and
denominator steps differ. The Cayley continuation supplies the integer-order
case; its Gamma-density argument extends without change to real outer order
$a\ge1$. It complements the sharper one-sided bounds above.

\begin{theorem}[Affine signed-moment bound]\label{affine:thm:moments}
Let $r,d$ be positive integers, $a\ge1$ real and $b>0$ real. Set
\[
 A_k=\frac{H_{rk}^{(b)}}{(dk+1)^a},\qquad
 m_1=1,\qquad m_a=\frac{(a-1)^{a-1}e^{-(a-1)}}{\Gamma(a)}\quad(a>1).
\]
There is a finite signed measure $\nu$ on $[0,1]$, with no atom at one,
such that $A_k=\int x^k\,d\nu(x)$ and
\[
 \|\nu\|_{\rm TV}\le\frac{2r}{d}m_a b\zeta(b+1).
\]
For $b>1$ one also has $\|\nu\|_{\rm TV}\le2\zeta(b)$.
Consequently the ordinary alternating series converges, and its $N$-term
Euler transform $E_N$ satisfies
\[
 \left|\sum_{k\ge0}(-1)^kA_k-E_N\right|\le C2^{-N},
\]
where $C$ may be either applicable bound. For integer $b$, valid rational
choices are $C=10r/(3d)$ at $b=1$ and $C=10/3$ at $b\ge2$.
\end{theorem}
\begin{proof}
Let $Y$ have Gamma density $y^{a-1}e^{-y}/\Gamma(a)$, and let $\mu$ be
the law of $X=e^{-dY}$. Then $\int x^k\,d\mu=(dk+1)^{-a}$.
In the logarithmic coordinate $s=-\log x$, its density is
\[
 f(s)=\frac{s^{a-1}e^{-s/d}}{d^a\Gamma(a)}\,\mathbf1_{s>0}.
\]
For every real $a\ge1$ its maximum is $m_a/d$. Extended by zero to the
negative axis, the total variation of its distributional derivative is
$2m_a/d$: the rise and fall give this for $a>1$, while at $a=1$ the
jump at zero and the decreasing exponential each contribute $1/d$.
The $L^1$ distance from a bounded-variation density to its translate by
$h\ge0$ is at most $h$ times this derivative variation. Indeed, integrate
$f(s-h)-f(s)=-\int_{s-h}^s f'(v)\,dv$ and use Fubini; smooth convolution
proves the same bound for the jump case. Thus, writing $D_\rho\mu$ for
pushforward by $x\mapsto\rho x$,
\[
 \|\mu-D_{e^{-rt}}\mu\|_{\rm TV}
 \le\min\left(2,\frac{2m_ar}{d}t\right).
\]
The generalized harmonic integral gives the measure
\[
 \nu=\frac1{\Gamma(b)}\int_0^\infty
   \frac{t^{b-1}}{e^t-1}(\mu-D_{e^{-rt}}\mu)\,dt.
\]
This integral converges in total variation: the translation bound supplies
an additional factor $t$ at zero, and exponential decay controls infinity.
Taking moments gives $A_k$, since the moments of the measure difference are
$(dk+1)^{-a}(1-e^{-rkt})$. Its total-variation bounds follow from
\[
 \frac1{\Gamma(b)}\int_0^\infty\frac{t^b}{e^t-1}\,dt
 =b\zeta(b+1),\qquad
 \frac1{\Gamma(b)}\int_0^\infty\frac{t^{b-1}}{e^t-1}\,dt
 =\zeta(b)\quad(b>1).
\]
Each constituent has no atom at one, and variation domination passes that
property to $\nu$. The finite alternating kernel is bounded by two and
converges for $x<1$ to $(1+x)^{-1}$, so dominated convergence proves
ordinary convergence. The exact remainder kernel in the Euler expansion is
$2^{-N}(1-x)^N/(1+x)$, bounded by $2^{-N}$; integrating against $\nu$
proves the stated error. Finally $m_a\le1$, by bounding the maximum of the
unit-scale Gamma density, and $\zeta(2)<5/3$. For the displayed rational
bound at $b=1$ one may also use only the integer $a$ needed by the rational
evaluator: then $m_a\le1$ follows directly from
$e^{a-1}\ge(a-1)^{a-1}/(a-1)!$. For the real-$a$ inequality, write the
Gamma density as the convolution of an exponential density and a
Gamma$(a-1,1)$ probability density when $a>1$; its supremum is at most one.
For integer $b\ge2$, $\zeta(b)\le\zeta(2)$ gives $10/3$.
\end{proof}
The real-order theorem does not make its centers rational at noninteger
indices. The finite rational implementation uses integer $a,b,r,d$ and
binomial-tail Euler weights. No complete monotonicity of $A_k$ is assumed:
the representing measure is signed and has total mass zero, since $A_0=0$.


\section{An exact kernel and the sharp Euler error}
\label{signed:sec:sharp-error}

\subsection{Separating the polynomial from the exponentially small part}
For this section $a,b$ are positive integers, $w=a+b$, and $T=-\log u$.

\begin{theorem}[Explicit density]\label{signed:thm:explicit-kernel}
Define the polynomial
\begin{align}\label{signed:eq:P-ab}
 P_{a,b}(T)={}&-\frac{T^{w-1}}{(w-1)!}\notag\\
 &+\sum_{\mu=b+1}^{w-1}(-1)^{\mu-b-1}\binom{\mu-1}{b-1}
       \zeta(\mu)\frac{T^{w-1-\mu}}{(w-1-\mu)!}.
\end{align}
Then, for $T>0$,
\begin{align}\label{signed:eq:explicit-kernel}
 \kappa_{a,b}(e^{-T})=P_{a,b}(T)
 +(-1)^{a-1}\sum_{j=0}^{b-1}\binom{w-j-2}{a-1}
       \frac{T^j}{j!}\Li_{w-1-j}(e^{-T}).
\end{align}
Consequently,
\begin{equation}\label{signed:eq:kernel-asymptotic}
 \kappa_{a,b}(u)=P_{a,b}(-\log u)
       +O\!\left(u(1+|\log u|^{b-1})\right)\quad(u\downarrow0).
\end{equation}
\end{theorem}
\begin{proof}
For $a=1$, expand $(e^v-1)^{-1}=\sum_{n\ge1}e^{-nv}$ in \eqref{signed:eq:kappa-base}. The elementary incomplete-gamma integral for integer $b$ gives
\[
 \frac1{(b-1)!}\int_T^\infty\frac{v^{b-1}}{e^v-1}\dd v
 =\sum_{j=0}^{b-1}\frac{T^j}{j!}\Li_{b-j}(e^{-T}),
\]
which is exactly the asserted base case.

Write $K_{a,b}(T)=\kappa_{a,b}(e^{-T})$. The recursion is
$K_{a,b}(T)=\int_0^T K_{a-1,b}(s)\dd s$. Differentiate the proposed right side of \eqref{signed:eq:explicit-kernel}, using
$\frac{\mathrm d}{\mathrm dT}\Li_s(e^{-T})=-\Li_{s-1}(e^{-T})$ and Pascal's identity. The polylogarithmic coefficients reduce to those for $(a-1,b)$; the derivative of $P_{a,b}$ is $P_{a-1,b}$. When $a\ge2$, the constant terms at $T=0$ cancel because
\[
 \binom{w-2}{b-1}=\binom{w-2}{a-1}
\]
and their signs are opposite. The asserted formula therefore has the correct derivative and initial value. This proves it by induction. Finally each $\Li_s(e^{-T})$ in the finite sum is $O(e^{-T})$ as $T\to\infty$, proving \eqref{signed:eq:kernel-asymptotic}.
\end{proof}

\subsection{The complete logarithmic error polynomial}
The exact error is
\begin{equation}\label{signed:eq:exact-error}
 E_N(a,b)-g_{a,b}
 =-2^{-N}\int_0^1\kappa_{a,b}(u)
                    \frac{(1-u^2)^N}{1+u^2}\dd u.
\end{equation}
This follows by summing only the first $N$ terms of the geometric kernel in the proof of Theorem~\ref{signed:thm:euler-certificate}. Its localization near $u=0$ explains why the polynomial in Theorem~\ref{signed:thm:explicit-kernel} determines the error.

\begin{theorem}[Sharp error with all logarithmic terms]\label{signed:thm:error-polynomial}
Define the polynomial of degree $w-1$
\begin{align}\label{signed:eq:Q-ab}
 Q_{a,b}(L)
 &=-\int_0^\infty e^{-t^2}P_{a,b}(L-\log t)\dd t\notag\\
 &=-\sum_{j=0}^{w-1}\frac{(-1)^jP_{a,b}^{(j)}(L)}{j!\,2^{j+1}}
       \Gamma^{(j)}(1/2).
\end{align}
For fixed $a,b$, as $N\to\infty$,
\begin{align}\label{signed:eq:error-polynomial}
 E_N(a,b)-g_{a,b}=2^{-N}\bigg\{&N^{-1/2}Q_{a,b}\!\left(\frac12\log N\right)\notag\\
 &+O\!\left(N^{-1}(1+\log^{b-1}N)
       +N^{-3/2}(1+\log^{w-1}N)\right)\bigg\}.
\end{align}
In particular,
\begin{equation}\label{signed:eq:error-leading}
 E_N(a,b)-g_{a,b}\sim
 \frac{\sqrt\pi}{2^w(w-1)!}
 \frac{2^{-N}(\log N)^{w-1}}{\sqrt N}.
\end{equation}
The leading coefficient depends only on the total weight, not on its split into $a$ and $b$.
\end{theorem}
\begin{proof}
Choose a fixed $\varepsilon\in(0,1/2)$. The part of \eqref{signed:eq:exact-error} with $u\ge\varepsilon$ is bounded by
$2^{-N}(1-\varepsilon^2)^N\|\kappa_{a,b}\|_1$ and is negligible at the displayed scales. In the remaining part substitute $u=t/\sqrt N$ and use \eqref{signed:eq:kernel-asymptotic}. For $t\le\varepsilon\sqrt N$,
\[
 \left|\frac{(1-t^2/N)^N}{1+t^2/N}-e^{-t^2}\right|
 \le\frac{C}{N}e^{-t^2}(t^2+t^4).
\]
Indeed, $0\le-\log(1-x)-x\le Cx^2$ for $0\le x\le\varepsilon^2$, and the denominator changes the expression by at most $t^2/N$ times the Gaussian bound.

The polynomial contribution, including the Jacobian $N^{-1/2}$, is therefore
$N^{-1/2}Q_{a,b}(\tfrac12\log N)$ up to
$O(N^{-3/2}(1+\log^{w-1}N))$. Extending the $t$ integral to infinity changes it only by an exponentially small amount. The remainder in \eqref{signed:eq:kernel-asymptotic} contributes at most
$O(N^{-1}(1+\log^{b-1}N))$; integrability of $t e^{-t^2}|\log t|^j$ at both endpoints justifies this estimate.

To obtain the second line of \eqref{signed:eq:Q-ab}, expand the polynomial and use
\[
 \int_0^\infty e^{-t^2}(\log t)^j\dd t
 =2^{-j-1}\Gamma^{(j)}(1/2),
\]
which follows by differentiating $\tfrac12\Gamma(s/2)$ at $s=1$. Finally the leading term of $-P_{a,b}$ is $T^{w-1}/(w-1)!$. Its Gaussian integral is $\sqrt\pi/2$, and $L=(\log N)/2$. This gives \eqref{signed:eq:error-leading}; both error terms in \eqref{signed:eq:error-polynomial} are smaller than that leading term.
\end{proof}

For example,
\begin{equation}\label{signed:eq:Q11}
 Q_{1,1}(L)=\frac{\sqrt\pi}{2}
       \left(L+\frac{\gamma+2\log2}{2}\right),
\end{equation}
where $\gamma$ is Euler's constant. Keeping the lower logarithmic terms matters at practical truncation lengths. For $(a,b)=(1,1)$, the following ratios use the exact rational proxy $E_N-E_{2N}$; its distance from the true error is at most $2^{-2N}$.
\begin{center}
\begin{tabular}{@{}rrr@{}}
\toprule
$N$ & Ratio to the leading term & Ratio to the full $Q_{1,1}$ approximation\\
\midrule
64 &1.42496&0.967964\\
128&1.37730&0.980510\\
256&1.33791&0.988050\\
512&1.30501&0.992593\\
\bottomrule
\end{tabular}
\end{center}
The asymptotic expansion is not substituted for the certified bound in the implementation. Its role is to explain the observed convergence and suggest more efficient future remainder enclosures.

\begin{proposition}[Euler certificate for the odd-denominator harmonic family]\label{signed:prop:S-euler}
For integer $p\ge1$, let $f(n)=H_n/(2n+1)^p$, let $d_k=D^kf(0)$, and let $E_N^S=\sum_{k<N}d_k/2^{k+1}$. Then
\begin{equation}\label{signed:eq:S-euler}
 E_N^S-\frac{N+1}{3^p2^N}\le S_p\le E_N^S,
 \qquad S_p=\sum_{n\ge0}\frac{(-1)^nH_n}{(2n+1)^p}.
\end{equation}
\end{proposition}
\begin{proof}
Use $H_n=\int_0^1(1-t^n)/(1-t)\dd t$ and the Mellin integral for $(2n+1)^{-p}$. Finite differencing gives
\[
 d_k=\frac1{\Gamma(p)}\int_0^1(-\log u)^{p-1}
  \int_0^1\frac{(1-u^2)^k-(1-u^2t)^k}{1-t}\dd t\dd u.
\]
Thus $d_0=0$ and $d_k<0$ for $k\ge1$. The numerator has absolute value at most $ku^2(1-t)$, so $|d_k|\le k/3^p$. Summing the Euler kernels yields
\[
 \sum_{k\ge0}\frac{d_k}{2^{k+1}}
 =-\frac1{\Gamma(p)}\int_0^1
  \frac{(-\log u)^{p-1}\log(1+u^2)}{1+u^2}\dd u=S_p.
\]
The final equality follows from the harmonic generating function, first with an Abel factor. The bound on $d_k$ justifies summing and integrating, and
$\sum_{k\ge N}k/2^{k+1}=(N+1)2^{-N}$ proves \eqref{signed:eq:S-euler}.
\end{proof}

At $N=400$, combine Proposition~\ref{signed:prop:S-euler}, Theorem~\ref{signed:thm:euler-certificate}, and the independent depth-one intervals. The exact rational computation encloses the difference between the two sides of \eqref{gaussian:eq:S4-short} in
\begin{equation}\label{signed:eq:S4-difference}
 [-10^{-118},\,10^{-118}].
\end{equation}
Both sides begin
\[
 -0.01049345629733504345335675100620175196351061887732\ldots.
\]
This is certified proximity, not proof of equality. The file
\texttt{s4\_interval\_certificate.json} records the outward-rounded enclosures and the tail bound used.

\subsection{The precise relation system being tested}
We now specify a finite-dimensional linear question; it must not be confused with completeness of all polylogarithm relations. Let
\[
 X_{a,b}^{r,s}=\Imag D_{a,b}(\iu^r,\iu^s),\qquad r,s\in\mathbb Z/4\mathbb Z,
 \qquad a+b=w.
\]
Set $X^{r,s}=0$ when both $r,s$ are even, and impose conjugation
$X^{r,s}=-X^{-r,-s}$. Choose the lexicographically smaller of each conjugate pair. Remove the divergent imaginary coordinate $X_{1,w-1}^{0,1}$; its conjugate is already identified. This leaves $6w-7$ formal coordinates.

For every $p+q=w$ and colors $r,s$, take the imaginary parts of the finite product stuffle
\begin{equation}\label{signed:eq:full-stuffle}
 D_{p,q}(x,y)+D_{q,p}(y,x)=\Li_p(x)\Li_q(y)-\Li_w(xy)
\end{equation}
and product shuffle
\begin{align}\label{signed:eq:full-shuffle}
 &\sum_{j=0}^{p-1}\binom{q+j-1}{j}D_{q+j,p-j}(y,x/y)\notag\\
 &\quad+\sum_{j=0}^{q-1}\binom{p+j-1}{j}D_{p+j,q-j}(x,y/x)
 =\Li_p(x)\Li_q(y),
\end{align}
where $x=\iu^r$, $y=\iu^s$. If exactly one factor is $\Li_1(1)$, retain the stuffle-minus-shuffle cancellation row: the common divergent double coordinate cancels, leaving the right side $-\Imag\Li_w(xy)$. Such rows follow by subtracting before taking the regularized boundary limit. Zero rows are discarded; duplicate nonzero rows need not be discarded. This construction defines a rational coefficient matrix $A_w$ and its depth-one right side. It is precisely the model implemented in \texttt{code/full\_ds.py}.

These are only linear depth-two product rows in a fixed weight. They do not include every relation obtained from distribution, argument transformations, higher-depth elimination, or the full algebraic consequences of standard relations. Zhao's treatment \cite{gaussian:Zhao} provides broader context for these distinctions.

\subsection{An exact obstruction at weight five}
Let $u_{4,1}=\Imag D_{4,1}(\iu,-1)$. Splitting the inner harmonic sum into even and odd indices proves
\begin{equation}\label{signed:eq:S4-color}
 S_4=g_{4,1}+u_{4,1}.
\end{equation}
Indeed, for an odd outer index $2n+1$, the sum of the two inner coefficients is
$H_{2n}+\sum_{m=1}^{2n}(-1)^m/m=H_n$; even outer indices have zero imaginary part.
Consequently, the left-hand coefficient vector needed for \eqref{gauss:eq:S4-closed} is
\begin{equation}\label{signed:eq:S4-target}
 T=3g_{4,1}+3g_{3,2}+9g_{2,3}+7u_{4,1}.
\end{equation}

\begin{proposition}[A finite linear-module obstruction]\label{signed:prop:S4-obstruction}
The coefficient vector $T$ in \eqref{signed:eq:S4-target} is not in the rational row span of the explicitly defined weight-five matrix $A_5$. Therefore no rational linear combination of just those weight-five depth-two rows proves the $S_4$ identity.
\end{proposition}
\begin{proof}
The matrix has 92 rows and 23 columns. In the canonical coordinate order above, define a vector $v$ by the following nonzero entries, setting every other entry to zero:
\begin{center}
\begin{tabular}{@{}crcr@{}}
\toprule
$(a,b,r,s)$ & $v$ & $(a,b,r,s)$ & $v$\\
\midrule
$(1,4,1,0)$&$-1/12$ & $(2,3,0,1)$&$1/24$\\
$(2,3,1,0)$&$1/8$   & $(2,3,1,3)$&$-1/24$\\
$(3,2,0,1)$&$-1/8$  & $(3,2,1,0)$&$-1/24$\\
$(3,2,1,3)$&$-1/24$ & $(4,1,0,1)$&$1/12$\\
\bottomrule
\end{tabular}
\end{center}
Substitution in the binomial rows \eqref{signed:eq:full-stuffle}--\eqref{signed:eq:full-shuffle} gives $A_5v=0$, while $T\cdot v=1$. These are finite rational equalities; the complete integer matrix, coordinate order, target, and witness are included in
\texttt{s4\_linear\_module\_obstruction.json} and replayed by \texttt{verify.py}. Every row-span vector annihilates $v$, whereas $T$ does not. This proves the assertion.
\end{proof}
The proposition concerns only the specified depth-two product rows.
Theorem~\ref{gaussian:thm:S4} adds a convergent octahedral duality and lifted
relations through higher depths, outside this row vocabulary. The two
statements are therefore compatible: the obstruction identifies the
additional information an eventual successful proof had to supply.

% ProveIt research continuation, 9 October 2026 (Pacific time).
% Replaces only the final rank-conjecture subsection of 05-signed-kernels.tex
% at commit 9bc738d3be22b8586a24693f19fb2e2a50ecd1bf.
% Full article, exact code and certificates accompany this fragment.

\subsection{Exact ranks from two polynomial reflections}
\label{l4rank:sec:rank}
Retain exactly the coefficient matrix $A_w$ defined above, including the
stuffle-minus-shuffle boundary rows and excluding the one divergent imaginary
coordinate. No distribution or higher-depth relations are added. Let
$A_w^{\neg g}$ denote the submatrix with the Gaussian columns deleted.
The following theorem replaces the former rank conjecture.

\begin{theorem}[All-weight ranks]
\label{signed:conj:rank}\label{l4rank:thm:rank}
For every odd $w\ge3$,
\begin{equation}\label{signed:eq:rank-conjecture}
 \operatorname{rank} A_w=5w-6,\qquad
 \operatorname{rank} A_w^{\neg g}=\frac{9w-11}{2}.
\end{equation}
For every even $w\ge2$,
\begin{equation}\label{l4rank:eq:even-rank}
 \operatorname{rank} A_w=6w-7,\qquad
 \operatorname{rank} A_w^{\neg g}=5w-6.
\end{equation}
In particular, the Gaussian-supported row space has dimension $(w-1)/2$
in odd weight and $w-1$ in even weight. In odd weight it is exactly the span
of the same-point shuffle rows.
\end{theorem}

\begin{proof}
Put $N=w-2$. Use six degree-$N$ polynomials $A,B,C,D,E,F$ to encode
a homogeneous null vector. Their respective colors are $(0,1),(1,0),(1,1),(1,2),(1,3),(2,1)$. Coefficients of $X^{a-1}Y^{b-1}$ have $a+b=w$.
Complete the missing coefficient of $A$ by
$v_{1,w-1}^{0,1}=-v_{w-1,1}^{1,0}$. This satisfies the missing stuffle row;
the missing shuffle row follows from its retained difference with stuffle.
The omitted coordinate occurs in no other row: outer index one in a
shuffle sum forces one of its factors to be $\Li_1(1)$.
Thus completion and restriction identify the matrix kernel with the full
homogeneous polynomial solution space.

Writing $P^{\mathsf T}(X,Y)=P(Y,X)$, the stuffle equations are
\begin{equation}\label{l4rank:eq:st}
 A=-B^{\mathsf T},\quad C=-C^{\mathsf T},\quad
 D=-F^{\mathsf T},\quad E=E^{\mathsf T}.
\end{equation}
The shuffle equations are
\begin{align}
 E(X+Y,X)+A(X+Y,Y)&=0,\notag\\
 B(X+Y,X)+B(X+Y,Y)&=0,\notag\\
 C(X+Y,Y)-F(X+Y,X)&=0,\notag\\
 D(X+Y,Y)-D(X+Y,X)&=0.\label{l4rank:eq:sh}
\end{align}
For general colors $(r,s)$ the shuffle polynomial is
$P_{s,r-s}(X+Y,X)+P_{r,s-r}(X+Y,Y)$; its coefficient of
$X^{p-1}Y^{q-1}$ is the stated binomial row. The four equations above come
from product colors $(0,1),(1,1),(1,2),(1,3)$. Swapping the factors,
conjugating, and discarding even-even imaginary rows exhaust the other cases.

The first two equations give
\[
 E(X,Y)=B(X-Y,X),\qquad B(X,Y)+B(X,X-Y)=0.
\]
With $\varepsilon=(-1)^N$, this implies
$E(Y,X)=-\varepsilon E(X,Y)$. Since $E=E^{\mathsf T}$, the whole $A,B,E$
block vanishes if $N$ is even. If $N=2m-1$ is odd, its free polynomial $B$
is precisely an anti-invariant under $Y\mapsto X-Y$, with basis
\begin{equation}\label{l4rank:eq:B-basis}
 B_j=X^{N-2j-1}(2Y-X)^{2j+1},\qquad 0\le j<m.
\end{equation}
Indeed, $U=X,V=2Y-X$ changes this reflection into $V\mapsto -V$.

The remaining equations give
\[
 F(X,Y)=C(X,X-Y),\quad D(X,Y)=-C(Y,Y-X),\quad C=-C^{\mathsf T}.
\]
They imply $D(X,X-Y)=-\varepsilon D(X,Y)$, while the last shuffle asks
for equality. Hence the $C,D,F$ block is zero when $N$ is even. When
$N=2m-1$, it has the basis determined by
\begin{equation}\label{l4rank:eq:C-basis}
 C_j=X^jY^{N-j}-X^{N-j}Y^j,\qquad 0\le j<m.
\end{equation}
Conversely, the displayed substitutions satisfy every homogeneous equation.
The $B_j$ and $C_j$ are independent and are retained in full after deleting
the missing coefficient of $A$. They therefore give an integral
$\mathbb Q$-basis of the full nullspace.

For odd $w=2m+1$, the nullity is $2m=w-1$, proving
$\operatorname{rank}A_w=(6w-7)-(w-1)=5w-6$. A null vector of
$A_w^{\neg g}$ extends by zero on the Gaussian columns; this is exactly
the $C$ block, of dimension $m$. Its $5w-6$ columns thus have rank
$5w-6-m=(9w-11)/2$. For even $w$ all columns of $A_w$ are independent,
and so are all columns of its submatrix.

Projection of the row space onto the non-Gaussian columns gives the stated
supported-row dimensions by taking the difference of the two ranks. In odd
weight $w=2m+1$, the same-point shuffle coefficient on $g_{a,w-a}$ is
$\binom{a-1}{p-1}+\binom{a-1}{w-p-1}$, $1\le p\le m$.
Its first $m$ columns form a triangular Pascal matrix with diagonal one.
These $m$ independent Gaussian-supported rows therefore exhaust that space.
\end{proof}

\paragraph{Exact membership and the harmonic-sum obstruction.}
Let $K_w$ have the preceding integer null vectors as columns. A rational
coefficient row $t$ is in the specified row space exactly when $tK_w=0$.
Otherwise any column with nonzero pairing is a separating certificate.
The basis is constructible in $O(w^2)$ integer arithmetic operations and
has $O(w)$-bit coefficients; coefficient recurrences and bit bounds are
given in the accompanying article.

For odd $w=2m+1$, choose $B=0$, $C=Y^{2m-1}-X^{2m-1}$.
It vanishes on every Gaussian column and has $D(1,0)=1$.
Since $S_{2m}=g_{2m,1}+\operatorname{Im}D_{2m,1}(\mathrm i,-1)$,
it pairs to one with $S_{2m}$. Thus no rational linear combination of these
fixed-weight product rows reduces $S_{2m}$ to Gaussian doubles and the
associated single-value right sides. In the formal quotient, adjoining it
to the Gaussian span increases the dimension from $m$ to $m+1$.
For the weight-five target above, this witness gives $T\cdot v=7$.
This is a uniform proof-system obstruction, not a numerical independence
theorem and compatible with the already proved $S_4$ identity, whose certificate uses a larger vocabulary.

\subsubsection{Restoring the right sides: an explicit compiler}
\label{l4rank:sec:compiler}
Set $\ell_k(r)=\Li_k(\mathrm i^r)$, with the formal regularization
$\ell_1(0)=0$, and define known degree-$N$ polynomials
\begin{align*}
 P_{r,s}(X,Y)&=\sum_{p+q=w}\operatorname{Im}(\ell_p(r)\ell_q(s))
                    X^{p-1}Y^{q-1},\\
 H_N(X,Y)&=\sum_{j=0}^{N}X^jY^{N-j},\\
 S_{01}&=P_{0,1}-\beta(w)H_N,\qquad
 S_{12}=P_{1,2}+\beta(w)H_N.
\end{align*}
The same six letters now encode the actual imaginary double values,
with the artificial missing coefficient $[X^0Y^N]A=-g_{w-1,1}-\beta(w)$.
Put
\begin{align*}
 e(X,Y)&=P_{0,1}(Y,X-Y)-S_{01}(X,X-Y),\\
 d(X,Y)&=S_{12}(X,Y)+P_{1,2}(X,Y-X),\\
 b(X,Y)&=P_{1,1}(Y,X-Y),\\
 h(X,Y)&=P_{1,3}(Y,Y-X)-e(Y,Y-X)+e(Y-X,Y).
\end{align*}

\begin{theorem}[Even-weight polynomial solution]\label{l4rank:thm:compiler}
For even $w\ge2$ the full solution is
\begin{align*}
 B(X,Y)&=\tfrac12\bigl(b(X,Y)+h(X,Y)\bigr),\\
 C(X,Y)&=\tfrac12\bigl(P_{1,3}(X,-Y)+P_{1,1}(Y,X)
                    -d(X-Y,-Y)+d(X-Y,X)\bigr),\\
 A(X,Y)&=S_{01}(X,Y)-B(Y,X),\\
 E(X,Y)&=B(X-Y,X)+e(X,Y),\\
 F(X,Y)&=C(X,X-Y)-P_{1,2}(Y,X-Y),\\
 D(X,Y)&=-C(Y,Y-X)+d(X,Y).
\end{align*}
Discarding the artificial coefficient gives explicit single-value reductions
for every convergent imaginary double at fourth roots of unity.
\end{theorem}
\begin{proof}
The affine stuffle equations are
$A+B^{\mathsf T}=S_{01}$, $C+C^{\mathsf T}=P_{1,1}$,
$D+F^{\mathsf T}=S_{12}$, $E-E^{\mathsf T}=P_{1,3}$.
The four shuffle right sides in \eqref{l4rank:eq:sh} are now
$P_{0,1},P_{1,1},P_{1,2},P_{1,3}$, respectively.
Elimination gives the last four displayed formulas and
$B(X,Y)+B(X,X-Y)=b(X,Y)$.
Even homogeneity and the $E$ stuffle give the difference
$B(X,Y)-B(X,X-Y)=h(X,Y)$, proving the formula for $B$.
For $C$, the last shuffle is
$D(u,u-v)-D(u,v)=P_{1,3}(v,u-v)$.
Use
$C(u-v,-v)=C(v-u,v)=P_{1,1}(v-u,v)-C(v,v-u)$
and then $(X,Y)=(v,v-u)$ to get the asserted formula.
Actual values satisfy these equations, and the zero homogeneous kernel
makes the solution unique.
\end{proof}

\paragraph{Two mixed leading-index families.}
Let $c_p=\operatorname{Im}D_{p,1}(\mathrm i,\mathrm i)$ and
$e_p=\operatorname{Im}D_{p,1}(\mathrm i,-\mathrm i)$. For all $m\ge1$,
\begin{align}
 c_{2m-1}&=(m-1)\beta(2m)-\tfrac12\log2\,\beta(2m-1)
 -\sum_{k=1}^{m-1}(1-2^{1-2k})\zeta(2k)\beta(2m-2k),
 \label{l4rank:eq:c-family}\\
 e_{2m-1}&=m\beta(2m)-\tfrac12\log2\,\beta(2m-1)
 -\sum_{k=1}^{m-1}\zeta(2k)\beta(2m-2k).
 \label{l4rank:eq:e-family}
\end{align}
For a direct coefficient proof, evaluate $C(1,0)$ and $E(1,0)$ in the compiler.
The resulting equations are
\begin{align*}
 2C(1,0)&=P_{1,3}(1,0)+P_{1,1}(0,1)-d(1,0)+d(1,1),\\
 2E(1,0)&=P_{1,1}(1,0)+P_{1,3}(1,0)+e(1,0)+e(0,1).
\end{align*}
In both lines the first two terms are $-\log2\,\beta(w-1)$, and
\begin{align*}
 d(1,1)-d(1,0)&=P_{1,2}(1,1)-P_{1,2}(1,-1)+N\beta(w),\\
 e(1,0)+e(0,1)&=-P_{0,1}(1,1)+P_{0,1}(1,-1)+w\beta(w).
\end{align*}
Only even inner, respectively even outer, indices survive the differences.
Using $\ell_{2k}(2)=-(1-2^{1-2k})\zeta(2k)$ and
$\operatorname{Im}\ell_j(1)=\beta(j)$ proves both formulas.

\paragraph{Status and evidence.}
The rank theorem proves both individual ranks, rather than inferring them
from their difference. The previous word ``Equivalently'' should therefore
be read as ``In particular.'' Independent matrix construction, exact integer
kernel checks and modular lower-bound certificates cover every weight
$2$--$41$; 210 even-weight formulas through weight 12 satisfy all 816 product
rows, and both leading-index formulas are checked through weight 16.
These finite checks support the uniform proofs; they are not their substitute.
The known general parity theorem already gives the relevant reduction
phenomenon. The new contribution here is the explicit solution of this
specified coefficient module. This row-system theorem asserts no numerical period independence. The
$S_4$ identity is proved separately in Theorem~\ref{gaussian:thm:S4}. Historical conjectural receipts are
retained as provenance rather than silently rewritten.

